\resheading{项目履历}
 \begin{itemize}[leftmargin=*,label={}]
 \item \begin{tabular*}{\textwidth-5mm}{ll@{\extracolsep{\fill}}r}
    \textbf{AI网络异构联合编译}& \quad{自研芯片后端接入负责人} & 2025.01-至今 \\
  \end{tabular*}
  {\small
    \begin{itemize}
      \item{完成自研芯片后端接入异构编译框架，支持算子映射与指令自动生成。}
      \item{协同业务需求方与芯片后端完成gridsample算子规格确认、优化加速与软硬件协同设计。}

    \end{itemize}
  }
 \item \begin{tabular*}{\textwidth-5mm}{ll@{\extracolsep{\fill}}r}
    \textbf{自研芯片迭代开发}& \quad{编译器领域PL} & 2023.06-至今 \\
  \end{tabular*}
  {\small
    \begin{itemize}
      \item{Pocket3产品手板阶段，SoC带宽争抢严峻，芯片对外DDR访存量遭质疑；利用编译插桩技术，直接在产品板上统计得到对外load/store指令量化数据，解决争议。}
      \item{芯片转阶段benchmark算子性能未达标；应用PGO、LTO技术，以较少投入大幅提升算子性能20\%，达标新一代芯片性能优化承诺，为项目阶段推进保驾护航。}
      \item{优化完善自研intrinsic编程模型，为设计、验证提供稳定API抽象层，节省验证人力投入1人月。}
      \item{深入算子开发前线，提供架构memory相关编程优化建议；承担部分产品化工作。}
    \end{itemize}
  }
  
  \item \begin{tabular*}{\textwidth-5mm}{ll@{\extracolsep{\fill}}r}
    \textbf{自研芯片功耗建模}& \quad{预研总负责人} & 2022.07-2024.02 \\
  \end{tabular*}
  {\small
  \begin{itemize}
    \item{协同硬件套片组完成功耗测量工具规格确认；优化测试流程，开发自动化功耗测量板驱动。}
    \item{标准化功耗测量方法，建立自研芯片算子功耗数据库资产。}
    \item{负责自研芯片功耗拟合工作，提出使用Roofline Model拟合profile数据，以计算访存比为唯一输入，勾勒出自研处理器的性能功耗曲线，支持产品化团队快速评估功耗。}
  \end{itemize}
  }
  
    %\item \textbf{求真品诚，担当开发主力}
%    {\small
%      \begin{itemize}
%        \item{基于LLVM的编译器开发，指令选择，DAG优化，别名分析，活跃变量分析等…}
%        
%        \item{：组织变革后，并跑通demo。}
%      \end{itemize}
%    }
%    \item \textbf{响应积极，关注投入产出比}
%    {\small
%      \begin{itemize}
%        
%      \end{itemize}
%    }
%    \item \textbf{融会贯通，掌握芯片全景视角}

%    \item \textbf{多面手}
%    
 \end{itemize} 